\Needspace{6\baselineskip}
\begin{center}\hypertarget{native-258-article185}{}\pdfbookmark[2]{185.}{article-185}185.\end{center}
Argumentum hactenus pertractatum pluribus disquisitionibus analyticis
elegantibus occasionem dare posset, quibus tamen hic non immoramur, ne
nimis ab instituto nostro distrahamur. Eadem ratione expositionem artificiorum,
per quae calculus numericus ad algorithmum magis expeditum reduci
potest, ad aliam occasionem nobis reservare debemus. Unicam observationem
hic adiicere liceat. Quoties multitudo functionum seu aequationum
\SourceLeaf{259}propositarum considerabilis est, calculus ideo potissimum paullo molestior evadit, quod
coëfficientes per quos aequationes primitivae multiplicandae sunt ut $P$, $Q$, $R$,
$S$ etc. obtineantur, plerumque fractiones decimales parum commodas involvunt.
Si in hoc casu operae pretium non videtur, has multiplicationes adiumento
tabularum logarithmicarum quam accuratissime perficere, in plerisque casibus
sufficiet, horum multiplicatorum loco alios ad calculum commodiores adhibere,
qui ab illis parum differant. Haecce licentia errores sensibiles producere nequit,
eo tantummodo casu excepto, ubi mensura praecisionis in determinatione
incognitarum multo minor evadit, quam praecisio observationum primitivarum
fuerat.
